% BHU PYQ -- Manually transcribed & error-corrected
% Paper: BCF-211 : Corporate Accounting
% Exam: B.Com. (Hons.) FMM Semester III Examination, 2022-23

\documentclass[11pt,a4paper]{article}
\usepackage[a4paper,top=2.5cm,bottom=2.5cm,left=2.8cm,right=2.8cm,headheight=14pt]{geometry}
\usepackage{amsmath,amssymb}
\usepackage[T1]{fontenc}
\usepackage[utf8]{inputenc}
\usepackage{lmodern,microtype,fancyhdr,enumitem,hyperref,lastpage,parskip,booktabs}

\hypersetup{
    colorlinks=true,
    urlcolor=black,
    linkcolor=black,
    pdftitle={BCF-211: Corporate Accounting -- BHU B.Com. (Hons.) FMM Sem III 2022-23}
}

\pagestyle{fancy}
\fancyhf{}
\renewcommand{\headrulewidth}{0.4pt}
\renewcommand{\footrulewidth}{0.4pt}
\fancyhead[L]{\small   \textit{Banaras Hindu University}}
\fancyhead[R]{\small   \textit{BCF-211: Corporate Accounting}}
\fancyfoot[C]{\small Page  \thepage\ of \pageref{LastPage}}


\newlist{parts}{enumerate}{2}
\setlist[parts,1]{label=(\alph*),leftmargin=*,itemsep=5pt,topsep=3pt}
\setlist[parts,2]{label=(\roman*),leftmargin=*,itemsep=3pt,topsep=2pt}

\newcommand{\pts}[1]{\hfill   \textit{[#1]}}


\begin{document}


\begin{center}
  {\large\bfseries BANARAS HINDU UNIVERSITY}\\[2pt]
  {
\normalsize B.Com. (Hons.) FMM --- Semester III Examination, 2022-23}\\[6pt]
  \rule{0.8\linewidth}{0.5pt}\\[6pt]
  {\large\bfseries Paper No. BCF-211: Corporate Accounting}\\[4pt]
  {
\normalsize Time: 3 Hours\hfill Maximum Marks: 70}
\end{center}

\rule{\linewidth}{0.4pt}

\smallskip



\noindent   \textit{Instructions: Answer all questions. All questions carry equal marks (14 marks each).}

\bigskip




\noindent   \textbf{Question 1.} \\
Ayush Ltd. issued for public subscription 20,000 equity shares of Rs. 10 each at a premium of Rs. 2 per share, payable as under:

\begin{itemize}
  \item On application: Rs. 2 per share
  \item On allotment: Rs. 5 per share (including premium)
  \item On first call: Rs. 2 per share
  \item On second call: Rs. 3 per share
\end{itemize}
Applications were received for 30,000 shares. Allotments were made pro-rata to the applicants for 24,000 shares, the remaining applications being rejected and their money returned. Money over-paid on application was utilized towards sums due on allotment.

A, to whom 800 shares were allotted, failed to pay the allotment and calls money; and B, to whom 1,000 shares were allotted, failed to pay the two calls. All the shares were subsequently forfeited after the second call was made. All the forfeited shares were sold to X as fully paid-up at Rs. 8 per share.

Show the Journal Entries required to record the above transactions in the books of Ayush Ltd.\dots\pts{14}

\centerline{   \textbf{OR}}

What do you understand by underwriting, sub-underwriting and firm underwriting? What are the various provisions in Companies Act, 2013 regarding Underwriting Commission?\pts{14}

\medskip\rule{\linewidth}{0.2pt}




\noindent   \textbf{Question 2.} \\
R Ltd. was absorbed by P Ltd. as on 31st March, 2022. All the assets and liabilities of R Ltd. were taken over by P Ltd. The purchase consideration was agreed at Rs. 3,36,600 and was payable in fully paid equity shares of P Ltd. distributed to the equity shareholders of R Ltd. 

The following are the Balance Sheets of both the companies as on 31st March, 2022:


\begin{table}[h!]
  \centering
  \small
  
\begin{tabular}{lrr}
     oprule
   \textbf{Particulars} &   \textbf{P Ltd. (Rs.)} &   \textbf{R Ltd. (Rs.)} \\
    \midrule
   \textbf{I. Equities and Liabilities:} & & \\
    (1) Shareholders' Funds: & & \\
    \quad Equity shares of Rs. 10 each, fully paid-up & 7,50,000 & 2,00,000 \\
    \quad Reserves and Surplus: & & \\
    \quad\quad General Reserve & 1,50,000 & 50,000 \\
    \quad\quad Profit \& Loss Account & 20,502 & 12,900 \\
    (2) Non-current Liabilities: & & \\
    \quad Workmen Compensation Fund & 12,000 & 9,000 \\
    \quad Staff Provident Fund & 10,000 & 4,000 \\
    (3) Current Liabilities: & & \\
    \quad Sundry Creditors & 58,567 & 30,456 \\
    \quad Provision for Taxation & 12,000 & 5,000 \\
    \midrule
   \textbf{Total} &   \textbf{10,13,069} &   \textbf{3,11,356} \\
    \midrule
   \textbf{II. Assets:} & & \\
    (1) Non-current Assets: & & \\
    \quad Goodwill & 2,00,000 & 60,000 \\
    \quad Plant and Machinery & 3,12,000 & 1,00,000 \\
    (2) Current Assets: & & \\
    \quad Stock in Trade & 2,65,000 & 80,000 \\
    \quad Sundry Debtors & 2,21,200 & 56,000 \\
    \quad Prepaid Insurance & --- & 700 \\
    \quad Income-tax Refund Claim & --- & 6,000 \\
    \quad Cash in Hand & 869 & 356 \\
    \quad Cash at Bank & 14,000 & 8,300 \\
    \midrule
   \textbf{Total} &   \textbf{10,13,069} &   \textbf{3,11,356} \\
    ottomrule
  \end{tabular}
\end{table}

You are required to:

\begin{parts}
  \item Show necessary Journal Entries in the books of R Ltd.
  \item Show necessary Ledger Accounts in the books of P Ltd.
  \item Prepare the new Balance Sheet of P Ltd. after amalgamation.
\end{parts}\pts{14}

\centerline{   \textbf{OR}}

What do you mean by amalgamation in the nature of merger and in the nature of purchase? What is the difference between the two?\dots\pts{14}

\medskip\rule{\linewidth}{0.2pt}

\pagebreak




\noindent   \textbf{Question 3.} \\
The Balance Sheet of Adarsh Ltd. as on 31st March, 2022 is as under:


\begin{table}[h!]
  \centering
  \small
  
\begin{tabular}{lr}
     oprule
   \textbf{Particulars} &   \textbf{Amount (Rs.)} \\
    \midrule
   \textbf{I. Equities and Liabilities:} & \\
    Share Capital: & \\
    \quad 7\% Preference shares of Rs. 10 each & 16,000 \\
    \quad Equity shares of Rs. 10 each & 40,000 \\
    \quad Profit \& Loss Account (deficit) & (22,260) \\
    \quad Less: General Reserve & 18,140 \\
    \quad Net Deficit & (4,120) \\
    Non-current Liabilities: & \\
    \quad 7\% Debentures of Rs. 10 each & 8,000 \\
    Current Liabilities: & \\
    \quad Creditors & 10,980 \\
    \midrule
   \textbf{Total} &   \textbf{70,860} \\
    \midrule
   \textbf{II. Assets:} & \\
    Non-current Assets: & \\
    \quad Fixed Assets & 18,680 \\
    Current Assets: & \\
    \quad Stock & 32,500 \\
    \quad Debtors & 18,700 \\
    \quad Cash & 980 \\
    \midrule
   \textbf{Total} &   \textbf{70,860} \\
    ottomrule
  \end{tabular}
\end{table}

A scheme of reconstruction duly approved by the court and implemented by the company is as under:

\begin{enumerate}[label=(\roman*)]
  \item Each Rs. 10, 7\% debentures be exchanged for one Rs. 5, 10\% debenture; one new 10\% preference share of Rs. 2.50 and one new equity share of Rs. 2.50 each.
  \item The preference shares (7\%) are to be reduced to Rs. 3.75 of which Rs. 1.75 will be represented by new equity shares and Rs. 2 by new 10\% preference shares.
  \item Equity shares are to be reduced to Rs. 2.50 each.
  \item Both preference shares and equity shares are to be consolidated into shares of Rs. 10 each.
  \item The Capital Reduction Account, after writing-off losses, is to be used in writing-off fixed assets and stocks proportionately.
\end{enumerate}

Give Journal Entries and prepare the Revised Balance Sheet of the Company.\pts{14}

\centerline{   \textbf{OR}}

What is Capital Reduction Account? When, why and how is it prepared?\pts{14}

\medskip\rule{\linewidth}{0.2pt}

\pagebreak




\noindent   \textbf{Question 4.} \\
Following was the position of Rudra Ltd. on 31st March, 2022:

\begin{itemize}
  \item Property: Rs. 90,000; Machinery: Rs. 30,000; Stock: Rs. 10,000.
  \item Profit \& Loss Account (Dr): Rs. 5,000.
  \item 5\% Debentures: Rs. 50,000; Loan on mortgage: Rs. 10,000.
  \item Unsecured creditors: Rs. 7,000; Income-tax (2020-21): Rs. 2,000.
  \item 4,000, 6\% Preference shares of Rs. 10 each; 2,600 Equity shares of Rs. 10 each.
\end{itemize}
Debentureholders have a floating charge on all the assets of the company. A receiver has been appointed by the debentureholders. There is voluntary liquidation of the company and a liquidator has been appointed. 

After taking into consideration the following information, prepare the Receipts \& Payments Account of the Receiver and the Liquidator's Final Statement of Account:

\begin{enumerate}
  \item Machinery is lodged as a security against mortgage debt. Machinery realized Rs. 25,000.
  \item Assets worth Rs. 80,000 were placed in the charge of the receiver who realized Rs. 78,000 on these assets. His remuneration is Rs. 400 and his cost is Rs. 200.
  \item Cost of liquidation is Rs. 800; Liquidator's remuneration is Rs. 500.
  \item Liquidator received Rs. 9,500 from the sale of remaining property.
  \item Legal expenses in liquidation Rs. 100.
\end{enumerate}\pts{14}

\centerline{   \textbf{OR}}

Define Holding Company and Subsidiary Company. What is Consolidated Balance Sheet and how is it prepared?\pts{14}

\medskip\rule{\linewidth}{0.2pt}




\noindent   \textbf{Question 5.} \\
Prepare a Statement of Affairs of a Company as per Companies Act, 2013.\pts{14}

\centerline{   \textbf{OR}}

Explain the following:\pts{7+7}

\begin{parts}
  \item Minority interest of shareholders and its calculation
  \item Cost of control
\end{parts}

\vfill
\rule{\linewidth}{0.4pt}

\begin{center}\small
  \href{https://bhu-exam.vercel.app/}{Click here to visit our website}\\[4pt]
  \href{https://me-aryan.vercel.app/}{Click here to visit our contributor}\\[4pt]
  \href{https://quashan-me.vercel.app/}{Click here to visit our contributor}
\end{center}

\end{document}
